%% kapa: exam=1 round=31 year=2020 period=1 form=A subject=real_estate subject_ko=부동산학원론 seq=23 number=103
%% source: Q-Net 감정평가사 기출문제 (cst00309) artlSeq=5207219 file=extracted/1008/5207219/A형/A형/2020년도 제31회 감정평가사 1차 시험_1교시_A형.pdf
%% conversion: kordoc 4.24.51 -> tools/gen1.py (auto); review=proofread-sonnet; flags=none
\begin{question}{103}
프로젝트 금융의 특징에 관한 설명으로 옳지 않은 것은?
\begin{choices}
  \item 사업자체의 현금흐름을 근거로 자금을 조달하고, 원리금 상환도 해당 사업에서 발생하는 현금흐름에 근거한다.
  \item 사업주의 입장에서는 비소구(non-recourse) 또는 제한적 소구(limited-recourse) 방식이므로 상환 의무가 제한되는 장점이 있다.
  \item 금융기관의 입장에서는 부외금융(off-balance sheet financing)에 의해 채무수용능력이 커지는 장점이 있다.
  \item 금융기관의 입장에서는 금리와 수수료 수준이 높아 일반적인 기업금융보다 높은 수익을 얻을 수 있는 장점이 있다.
  \item 복잡한 계약에 따른 사업의 지연과 이해당사자 간의 조정의 어려움은 사업주와 금융기관 모두의 입장에서 단점으로 작용한다.
\end{choices}
\end{question}
\answer{3}
